\section{Введение}

Длины волн обозначаются $\lambda^{\text{т}}_{\text{в}}$, $\lambda^{\text{т}}_{\text{с}}$, $\lambda^{\text{п}}_{\text{в}}$ и $\lambda^{\text{п}}_{\text{с}}$. Верхний индекс указывает на теоретическую или экспериментальную оценку, нижний~--- на волновод (в) или свободное пространство (с). Индекс~2 дополнительно обозначает оценку по измерениям второго задания.
\subsection*{Цель} \addcontentsline{toc}{subsection}{Цель}

Нахождение коэффициента направленного действия (КНД) пирамидальной рупорной антенны зеркальным методом (методом Парселла).

\subsection*{Задачи} \addcontentsline{toc}{subsection}{Задачи}

1. Помещая перед раскрывом рупора щит с поглощающим покрытием ($\Gamma \approx 0$), снять зависимость $|E|^2(x)$ вдоль волновода. Определить $\Gamma_{\text{к}}$ и оценить КПД системы по рассогласованию. По длине волны в волноводе ${\lambda^{\text{п}}_{\text{в}}}$ найти длину волны в свободном пространстве ${\lambda^{\text{п}}_{\text{с}}}$.

2. Выбрать положение зонда, при котором $\cos(2hx + \varphi_{\text{к}}) = 0$, и зафиксировать его. Убрав поглощающий щит и меняя положение антенны, снять $|E|^2(\Delta X)$. Определить $\Gamma$ и по нему КНД антенны. Убедиться в малости $\Gamma_{\text{к}}^2$, $\Gamma^2$, $\Gamma_{\text{к}}\Gamma$. Выяснить, как по полученной зависимости определить ${\lambda^{\text{п}}_{\text{с}}}$.

3. Зафиксировать положение антенны относительно щита и измерить $\kappa = E_{\min}/E_{\max}$ в зависимости от $X + \Delta X$. Снять зависимость
\[
\tilde{\Gamma}(\Delta X) = \frac{1 - \kappa}{1 + \kappa} = \frac{1 - \sqrt{K}}{1 + \sqrt{K}}.
\]
Показать, что $\tilde{\Gamma}_{\max} \approx \Gamma + \Gamma_{\text{к}}$, и, используя найденное $\Gamma_{\text{к}}$, определить $\Gamma$ и КНД. Сопоставить результаты с п.~2 и оценить погрешности обоих методов.

4. По принципу Гюйгенса в кирхгофовском приближении рассчитать теоретический КНД рупора, используя упрощающее предположение о синфазности поля на раскрыве. Обсудить, является ли поле синфазным в действительности и как неоднородность фазы влияет на КНД.

5. Сравнить экспериментальные и теоретические результаты, объяснить причины расхождений.

\subsection*{Приборы и оборудование} \addcontentsline{toc}{subsection}{Приборы и оборудование}
\begin{itemize}
	\item Генератор;
	\item Измерительные линейки;
	\item Амперметр квадратичного детектора;
	\item Рупорная антенна (апертура $l_1 \times l_2 = 9.2 \times 13.4~\text{см}$);
	\item Поглощающий щит;
	\item Отражающий щит.
\end{itemize}
